% Transcription — fonds Alexandre Grothendieck, shelfmark 156-4, batch 4 (pages 61-80).
% Established from the facsimile archives/batches/156-4/batch-04.pdf (a copy of
% 156-4.pdf fetched through the project's relay; no cover sheet in the batch, so
% PDF page k = folder page 60+k), per the skill transcribe-grothendieck. The
% folder is eighty-eight pages in five batches; this is the fourth. Batches 1,
% 2, 3 and 5 were read for notation. Page 61 continues batch 3's page 60,
% which reconstructs (𝔉, ≤) from (𝓜, ≤, R); the axioms C 1-C 3 it cites are
% batch 3's (pp. 53-58) and are not restated. This batch carries C4-C7 of
% that run, which batch 5 cites.
%
% Pass: Opus 5.5 (claude-opus-5-5), 2026-09-26 — first pass, unchecked against the pages by a human.
%
% Pages skipped: none. Pages summarised: none. All twenty leaves are his
% mathematics, fast drafting in ink. Cancelled blocks (long oblique strokes)
% are transcribed and noted once each: page 62 (Corollaire 2 and the
% paragraph after it), page 70 (lower half: an upside-down first start of a
% page numbered « 55 », on the sets A_X and irreducible objects), page 72
% (the conditions a), b)), page 75 (a first attempt on induced refinement),
% page 76 (the proof), page 78 (« C 8 »), page 80 (the proof at the foot).
%
% His pagination 61-80 runs at the top of the leaves and coincides with the
% archivists' (page 70's number is written over another figure). One date in
% his hand: « 14 juin », underlined, in the left margin of page 80, where a
% new paragraph begins; it is set as a subsection heading with a note, as
% batches 2 and 3 set « 12 juin » and « 13 juin ». No other date on
% these pages. His numbered runs: the axioms C1, C2, C3 (cited, p. 61; from
% batch 3), C4 (p. 64), C4 bis (« Axiome renforcé », p. 65), C5 (p. 69),
% C6 with (a), (b) circled and a marginal (c) (p. 71), C7_0 and C7 (p. 72),
% « C 8 » stated and cancelled (p. 78); Corollaire 1, 2 (p. 62), Corollaire
% 1, Cor. 2, Cor 3 (p. 64), « Définition 1 » (p. 71), Lemme 1 (p. 77) and
% Lemme 2 (p. 78). His cross-references: « cf. p. 57 » (p. 62), « cf p. 67 »
% (p. 70), « Cont 4 (p. 35) » and « p. 47 » (p. 75), « page 51 » (p. 78),
% « prop. p. 78 » and « C4 p. 47 » (p. 80, the 78 overwritten).
%
% What the batch is about. The combinatorics of a contrée seen as a set 𝔉 of
% figures with the sub-figure order ≤ (axioms C1-C3 of batch 3), its
% irreducible elements the multistrates 𝓜, and a compatibility relation R
% (the « quatrième mouture » of the contrée opened on page 52, batch 3).
% Page 61: for « finite » figures the data (𝔉, ≤) amount to (𝓜, ≤, R),
% figures being the finite closed parts of 𝓜 whose elements are pairwise
% compatible. Pages 61-65: connected and 0-connected figures, connected
% components of a figure and of 𝓜 = ∐ 𝓜_i, the map 𝔉 → ∏ 𝔉_i, axioms C4
% and C4 bis, π_0. Pages 65-68: cosupport and support of a set of figures,
% saturated parts, the set Σ of « supports virtuels » with arbitrary Inf and
% Sup, which he presumes a Boolean algebra, « ultra-stonienne ». Pages 69-73:
% the second order ≼ (subdivision), « contrées modérées », axioms C5 (a
% subdivision has the same disjunctions), C6 (the induced subdivision of a
% sub-figure), C7_0 and C7, sub-contrées (Définition 1, p. 71), the relation
% « G raffine F », G ≪ F iff G ≤ F' ≼ F, and the proof that it is an order,
% reduced to the Lemme F ≼ K ≤ F ⟹ K = F. Pages 74-75: that lemma is not
% formal in (𝔉, ≤); an example with 𝔉 all closed parts of 𝓜; a missing
% axiom. Pages 76-80: the refinement of a sub-figure induced by a refinement
% (Prop., p. 76), Inf for ≪, Lemmes 1-2 (pp. 77-78), which he abandons
% (« J'abandonne ici ! », p. 78, reading uncertain); a Proposition (p. 78)
% that for X ≪ F there is a least strate Y of F with X ≪ Y, proved on page
% 79; the increasing map φ : F̄ → Ḡ (p. 80); and, on 14 June, the notation
% F̃ for the set of X ≪ F and a corollary G' ≼ G ≼ F ⟹ G̃' ∩ F̃ ⊂ G̃, whose
% proof page 81 (batch 5) continues.
%
% Notation fixed once, aligned with batches 2 and 5: \mathcal{M}
% (multistrates), R (compatibility, as batch 3), \mathfrak{F} (figures),
% \mathfrak{P} (power set), \preccurlyeq for his curved second order ≼
% (subdivision), \ll for « raffine » (he writes two chevrons, sometimes
% three), \overline for his bar (the set of strates of a figure),
% \widetilde for his new F̃ (p. 80), \amalg / \coprod for his disjoint
% union. His stacked inequalities (vertical ≤, ≼ between rows) are set as
% arrays with the relation placed between the two rows, read from the upper
% term to the lower. His circled (a), (b), (c) are set in parentheses with a
% note. His underlining is \emph.
%
% Hard pages: 61 (an ink blot and a long oblique note across the text), 62,
% 68 (a margin note running the height of the page), 71, 75 and 80 (the
% interlinear insertions). About 195 \ill{} and 19 \uncertain{}; most of the
% loss is connective prose, as usual on his fast pages. Readings flagged and
% not settled: the exponent of π_0(𝔉) (p. 65); « compatibles » /
% « comparables » (p. 69), where the page may say either; « J'abandonne
% ici » (p. 78); « Supposons X ≤ F » at the foot of page 80 (read « ∈ » at
% first; batch 5 reads ≤); the two relations in « G' ≼ G ≼ F » (p. 80),
% both inked over. Variances on the page, not repaired: « sous-figure
% connexe de G » for « de F » (p. 61); « 𝔉_0 de 𝓜 » (p. 67); « une plus
% grande sous-figure F' de G » for G' (p. 76); « H ≤ G » and « Y ≪ H » in the
% cancelled proof of page 76; in C6 (p. 71) the statement names H the
% figure its square calls G'.
%
\documentclass[11pt,a4paper]{article}
\input{../preamble/grothendieck}

\folder{156-4}
\batch{4}
\pages{61}{80}
\dating{1986}
\watermark{Édition de démonstration}
\foldertitle{[Chapitre] IV. Analysis situs (première mouture) : notes manuscrites (10/06/1986)}

\begin{document}

\page{61}
\struck{Soit} Ainsi, les systèmes $(\mathfrak{F}, \leq)$ satisfaisant C1, C2, C3 et tels que toute $F \in \mathfrak{F}$ soit « finie » (i.e. l'ens. $\mathfrak{F}_{\leq F}$ des sous-figures est fini, ou encore, l'ens. $\mathcal{M}_{\leq F} = F$ \struck{\ill{}} de \struck{\ill{}} ses strates est fini) équivalent à : $(\mathcal{M}, \leq, R)$ \struck{\ill{}} où $\leq$ relation d'ordre, $R$ relation réfl. et sym., satisfaisant
\begin{enumerate}
  \item[a)] les $\mathcal{M}_{\leq X}$ finis ($X \in \mathcal{M}$)
  \item[b)] si $(X,Y) \in R$ et $X' \leq X$, $Y' \leq Y$, on a $(X',Y') \in R$ ;
\end{enumerate}
où les figures sont \ill{} parties \add{$A$} \ill{} les parties \emph{fermées}, \emph{finies} de $\mathcal{M}$, telles que \ill{} $X, Y \in A$, on ait $(X,Y) \in R$.
\note{les conditions C1, C2, C3 sont celles des pages précédentes (page 58 pour C3) ; une longue note oblique, écrite en travers du bas de cet alinéa et tachée d'encre, n'est lue qu'en fragments (« \ill{} la relation de compatibilité \ill{} les strates \ill{} En fait, \ill{} ») ; au-dessus de « fermées » un mot inséré n'est pas lu}

\emph{Composantes connexes} \quad Une figure est \emph{connexe} si \struck{\ill{}} (\add{avec $F'$, $F''$ \emph{disjointes} i.e. $F' \cap F'' = \emptyset$})
\[
  F = F' \cup F'' \Longrightarrow F' = F \ \text{ou}\ F'' = F .
\]
\note{les seconds membres $F$ sont écrits en surcharge sur d'autres lettres}
\marginal{NB \ill{} $F$ connexe \ill{}}
Une figure est dite \emph{0-connexe} si elle est connexe et $\neq \emptyset_{\mathfrak{F}}$. Une sous-figure $F_0$ de $F$ est dite \emph{composante connexe} de $F$ si elle est 0-connexe et $F = F_0 \amalg G$, \struck{$G$ \ill{}} $F_1$ disjoint de $F_0$.

b) \struck{$F_0$ est une composante}

\struck{Lemme} \struck{NB} Soit $F = F' \amalg F''$, \struck{avec $F' \cap F'' = \emptyset$}, et soit $G$ une sous-figure connexe de $G$. Alors $G$ est contenue dans $F'$ ou $F''$.
\note{« de $G$ » pour « de $F$ » sur la page}

\marginal{On a $F = F' \amalg F''$ \ill{} chaque figure \ill{} disjointe \ill{} $F'$, $F''$ \ill{} (NB si \ill{}) \ill{}}
\note{note oblique de la marge gauche, en bas de page, lue en partie seulement}

\page{62}
En effet, $G = G' \amalg G''$ avec $G' \subset F'$, $G'' \subset F''$ \struck{$G'$ et $G''$ disjoints} (cf. p. 57) donc on a $G' = G$ ou $G'' = G$, donc $G \subset F'$ ou $G \subset F''$ cqfd.

\emph{Corollaire 1} \quad \struck{Deux} Soient $F_0$, $F_1$ deux comp. connexes de $F$. Alors ou $F_0 = F_1$, ou $F_0$ et $F_1$ sont disjoints, et \struck{\ill{}} \add{\uncertain{dans ce cas}} $F = F_0 \amalg F_1$ \struck{$\amalg H$}.
\note{« Corollaire 1 » est souligné et une accolade verticale borde l'énoncé à gauche}

On écrit $F = F_0 \amalg G$, et on \ill{} que $G$ \ill{}

\emph{Corollaire 2} \quad Soit $(F_i)_{i \in I}$ une famille de composantes connexes deux à deux distinctes. Alors \struck{$F_0 \amalg F_1$} \ill{} figures \ill{} par les $F_i$ est \ill{} réunion disjointe $\coprod F_i \subset F$, et si la famille est finie (p. ex. $F$ \struck{fini} figure finie) alors
\[
  F = \Bigl(\coprod F_i\Bigr) \amalg H .
\]

Je n'ai pas pu déduire que \struck{\ill{}} figure \ill{} soit \struck{\ill{}} \ill{} comp. connexes, \struck{Mais \ill{} \ill{} des \ill{}, \ill{} \ill{} et \ill{} \ill{} pour les \ill{} des sous-figures « \ill{} »} figures finies. Les
\note{le Corollaire 2 et l'alinéa qui le suit sont barrés de cinq longs traits obliques}

\emph{Proposition} \quad Soit $F$ une figure, $A = \overline{F} \subset \mathcal{M}$ l'ens. de ses strates. \struck{Alors l'application} Pour qu'une sous-figure \struck{$F_0$} $F_0$ de $F$ soit une comp. conn. de $F$, il faut et il suffit que $\overline{F_0}$ soit une composante connexe de l'ens. ordonné $A$ (pour l'ordre induit par $\leq$). On trouve ainsi une corr. biunivoque entre l'ens. des composantes

\page{63}
de $F$, et celles de $A$. De plus, $F$ est somme directe disjointe de ses comp. connexes, i.e. $F = \operatorname{Sup}_i F_i$, les comp. connexes $F_i$ deux à deux disjointes…
\marginal{NB \ill{} C3 \ill{}}
\note{note oblique de la marge gauche, en haut de page, lue en fragments}

\emph{Composantes connexes de \uncertain{la} contrée $\mathfrak{F}$}

Soit $\mathcal{M} = \coprod_{i \in I} \mathcal{M}_i$ la décomposition de $\mathcal{M}$ en composantes connexes. Pour toute figure $F$, posons
\[
  F_i = \operatorname{Sup}_{X \in \mathcal{M}_i \cap \mathcal{M}_{\leq F}} X .
\]
Alors $F$ est la somme directe des $F_i$, \struck{\ill{}} si $F$ est de t.f. (\uncertain{ou} $I = \pi_0(\mathcal{M})$ fini!) alors tous les $F_i$ \uncertain{sauf} un nb fini \struck{\ill{}} $\emptyset$.
\marginal{\ill{} les composantes \ill{} $F_i$ \ill{} de $F$ sur $\mathcal{M}_i$}
\note{note oblique de la marge gauche, lue en partie seulement}

\struck{Inversement} Soit $\mathfrak{F}_i$ l'ens. des $F \in \mathfrak{F}$ telles que $F \subset \mathcal{M}_i$. Alors \struck{\ill{}} $\mathfrak{F}_i$ satisfait les conditions C1, C2, \struck{\ill{}} (et C3 aussi, si $\mathfrak{F}$ y satisfait), et $\mathcal{M}_i$ est l'ens. de ses objets irréductibles. \struck{L'application} Si $F \in \mathfrak{F}$, si $F_i$ sont ses composantes suivant les $\mathcal{M}_i$, \struck{$\mathfrak{F} \to \prod \mathfrak{F}_i$ \ill{} $F \mapsto (F_i)_{i \in I}$ \ill{} $F_i$ \ill{}} alors
\[
  \overline{F_i} = \overline{F} \cap \mathcal{M}_i , \qquad \overline{F} = \bigcup \overline{F_i} \quad (\text{réunion disj.}) .
\]
Pour $i$ fixé, l'application $\mathfrak{F} \to \mathfrak{F}_i$ commute aux Sup quelc. et aux Inf \uncertain{non vides}. L'application
\[
  \mathfrak{F} \to \prod \mathfrak{F}_i , \qquad F \mapsto (F_i)
\]
est injective, et la relation d'ordre de $\mathfrak{F}$ est induite par l'ordre produit

\page{64}
Inversement, donnons-nous, \struck{$F_i \in$} $\forall i \in I$, un $F_i \in \mathfrak{F}_i$, pouvons-nous \uncertain{affirmer} qu'il existe $F \in \mathfrak{F}$ telle que $\forall i$, $F_i$ soit la $i$-composante de $F$ ?

Déjà, il le faut \uncertain{sûrement} pour \struck{\ill{}} \ill{} tous les $F_i$ \uncertain{sauf} un nb fini soient \struck{$\emptyset$} vides. Donc on supposera donc

\emph{C4} \quad Si $X, Y \in \mathcal{M}$ appartiennent à deux comp. connexes disjointes de $\mathcal{M}$, alors $X$ et $Y$ sont compatibles.

C'est nécessaire et suffisant pour la \ill{} des \ill{}

\struck{Cor}

\emph{Corollaire 1} \quad L'image de $\mathfrak{F} \to \prod \mathfrak{F}_i$ contient les $(F_i)_{i \in I}$ tels que $F_i = \emptyset$ sauf pour \struck{\ill{}} un nb fini d'$i$.

\emph{Cor. 2} \quad \struck{Donc} $\mathfrak{F} \to \prod \mathfrak{F}_i$ est une bijection dans les \struck{deux} cas où $I$ est fini \struck{suivants : a) $I$ est fini ; b) toute figure est de t.f. \ill{} alors}
\note{« Cor. 2 » est écrit dans un cadre tracé à l'intérieur du Corollaire 1}

\emph{Cor 3} \quad Supposons que les figures \struck{\ill{}} $\in \mathfrak{F}$ soient de t.f. (resp. finies) \struck{\ill{}}. Alors il en est de même pour $\mathfrak{F}_i$, de plus $\mathfrak{F} \to \prod \mathfrak{F}_i$ a comme image l'ens. des $(F_i)$ tels que $F_i$ soit $\emptyset$ pour presque tout $i$.

\page{65}
Ceci résulte de \struck{C4 et} cor. 1 et du

\emph{Lemme} \quad Pour que $F$ soit de t.f. (resp. fini) il faut et il suffit que a) les $F_i$ le soient et b) tous les $F_i$ sauf un nb fini soient $\emptyset$.

\emph{Notations} \quad $\pi_0(F)$ (NB si $F \leq G$ on a $\pi_0(F) \to \pi_0(G)$ \ill{}, et transitivité), $\pi_0(\mathcal{M})$, $\pi_0^{\cdot}(\mathfrak{F})$. (Attention, le $\pi_0^{\cdot}(\mathfrak{F})$ n'est pas le $\pi_0(\mathfrak{F}, \leq)$, mais le $\pi_0$ \ill{} en tenant compte des « Sup ».) Si $\mathcal{C} = (\mathfrak{F}, \leq, \ldots)$, on peut aussi écrire $\pi_0(\mathcal{C})$.
\note{l'exposant de $\pi_0(\mathfrak{F})$, un point surmonté d'un signe biffé puis d'un « d » à sa seconde occurrence, est douteux ; sous $\mathcal{C}$ il écrit « contrée », le troisième terme du triplet, rendu ici « $\ldots$ », est un signe biffé non lu, sous lequel deux lignes ne se lisent qu'en partie (« \ill{} vues ici »)}

\emph{Axiome \uncertain{renforcé}}

\emph{C4 bis} \quad Pour tout $(F_i)_i \in \prod \mathfrak{F}_i$, \struck{\ill{}} $\exists\, F \in \mathfrak{F}$ telle que $F|\mathcal{M}_i = F_i$ pour $i \in I$.

\emph{Contrée résiduelle} \quad \emph{Support et complément}

Soit \struck{\ill{}} $\Phi \subset \mathfrak{F}$ \add{$H_i$} un ens. de figures. \struck{On posera} Soit $F$ une figure. Pour \struck{Compl.\ $\Phi$} que $F$ soit disjointe de tous les $H_i$, il faut et il suffit que ce soit le cas pour toutes les strates de $F$ (cf. conséquence de C3). \struck{Posons} Posons

\page{66}
\begin{align*}
  \mathfrak{F}' &= \operatorname{cosupp}_{\mathfrak{F}}(\Phi) = \lbrace F \in \mathfrak{F} \mid F \ \text{disjoint des}\ H_i \in \Phi \rbrace \quad (\subset \mathfrak{F}) \\
  \mathcal{M}' &= \operatorname{cosupp}_{\mathcal{M}}(\Phi) = \lbrace X \in \mathcal{M} \mid X \ \text{disjoint des}\ H_i \in \Phi \rbrace \\
  &= \mathcal{M} \cap \operatorname{cosupp}_{\mathfrak{F}}(\Phi) \subset \mathcal{M}
\end{align*}

\struck{Donc \ill{}, pour $F \in \mathfrak{F}$, posant $\mathfrak{F}' = \operatorname{cosupp}_{\mathfrak{F}}(\Phi)$, $\mathcal{M}' = \mathfrak{F}' \cap \mathcal{M} = \operatorname{cosupp}_{\mathcal{M}}(\mathfrak{F})$} \struck{$\Phi$} \struck{$F \in \operatorname{cosupp}$}

Donc pour $F \in \mathfrak{F}$
\[
  \left\lbrace
  \begin{array}{l}
    F \in \mathfrak{F}' \Longleftrightarrow \overline{F} \subset \mathcal{M}' \\
    \mathcal{M}' = \mathfrak{F}' \cap \mathcal{M}
  \end{array}
  \right.
\]
\marginal{ainsi, $\mathcal{M}'$ et $\mathfrak{F}'$ se déterminent mutuellement}
\marginal{NB $\mathfrak{F}'$, $\mathcal{M}'$ sont des parties fermées de $\mathfrak{F}$, $\mathcal{M}$. \ill{} $\mathfrak{F}'$ (stables par sup.)}
\note{la première note est en marge gauche, avec une flèche vers l'accolade ; la seconde en marge droite}

Soit d'autre part
\begin{align*}
  \mathfrak{F}'' &= \operatorname{cosupp}_{\mathfrak{F}}(\mathfrak{F}') \overset{\text{déf}}{=} \operatorname{supp}_{\mathfrak{F}}(\Phi) \\
  \mathcal{M}'' &= \operatorname{cosupp}_{\mathcal{M}}(\mathfrak{F}') \overset{\text{déf}}{=} \operatorname{supp}_{\mathcal{M}}(\Phi)
\end{align*}
On a donc $\mathfrak{F}'' \supset \Phi$, $\mathfrak{F}''$ stable par Sup. \struck{\ill{} $\Phi$} et par passage aux sous-figures. Soit $\overline{\Phi} = \bigcup_{H_i \in \Phi} \overline{H_i}$, alors
\begin{align*}
  \operatorname{cosupp}_{\mathfrak{F}}(\Phi) &= \operatorname{cosupp}_{\mathfrak{F}}(\mathcal{M}_{\overline{\Phi}}) \\
  \operatorname{cosupp}_{\mathcal{M}}(\Phi) &= \lbrace X \in \mathcal{M} \mid X \ \text{disjoint de tt}\ Y \in \mathcal{M}_{\overline{\Phi}} \rbrace
\end{align*}
donc on est ramené à des constructions dans $\mathcal{M}$. §
\note{« passage aux sous-figures » : lecture probable d'une fin de ligne écrite vite}

\emph{Partie saturée de $\mathfrak{F}$} \struck{$\mathfrak{F}_0$} : égale au cosupport (\uncertain{p.r. à} $\mathfrak{F}$) de son cosupport. Partie saturée \add{$\mathcal{M}_0$} de $\mathcal{M}$ : égale au cosupport de son cosupport (dans $\mathcal{M}$). Elles se correspondent mutuellement par
\[
  \mathcal{M}_0 = \mathfrak{F}_0 \cap \mathcal{M} , \qquad \mathfrak{F}_0 = \lbrace F \in \mathfrak{F} \mid F \subset \mathcal{M}_0 \rbrace .
\]

\page{67}
\struck{On a une notion de \ill{}}

Toute partie de $\mathfrak{F}$ (resp. $\mathcal{M}$) est contenue dans une partie saturée minimale (parmi celles la contenant), qui \struck{s'appelle le} n'est autre que le cosupport du cosupport, et s'appelle aussi son support. Les parties saturées sont celles qui sont égales au support, ou au cosupport, d'une partie. L'application cosupp définit une bijection de l'ens. des parties saturées sur \ill{}, renversant la relation d'inclusion.

On désigne par $\Sigma$ l'ens. des \emph{supports virtuels}, ou parties saturées \add{$\mathcal{M}_0$} de $\mathcal{M}$, ou \struck{de} $\mathfrak{F}_0$ de $\mathcal{M}$. \struck{\ill{}}

\struck{L'application de cosupport d}

On \ill{} $\Sigma$ à \struck{\ill{}} un \ill{} (comme partie de $\mathfrak{P}(\mathcal{M})$, ou de $\mathfrak{P}(\mathfrak{F})$) et on écrit, pour $\sigma \in \Sigma$,
\[
  \mathcal{M}_\sigma \subset \mathcal{M}, \qquad \mathfrak{F}_\sigma \subset \mathfrak{F} .
\]
Dans $\Sigma$, \struck{pour} la notion de cosupport est une anti\ill{} d'ens. ordonnés. Dans $\Sigma$, on a des Inf quelconques (c'est l'intersection des parties de $\mathcal{M}$, ou de $\mathfrak{F}$), et des Sup quelconques (c'est la partie saturée
\note{« $\mathfrak{F}_0$ de $\mathcal{M}$ » : lire « de $\mathfrak{F}$ » ; la lettre finale est surchargée}

\page{68}
engendrée par la réunion). Je \struck{dis} \add{présume} qu'on a en fait une algèbre de Boole (« ultra-stonienne », à cause des Sup quelc.)

\struck{Il faudra y revenir} \ill{} les supports s'identifient aux \ill{} \ill{} et \ill{} \ill{} espace, qui serait un espace \ill{} discontinu ultra-stonien. \ill{} correspondant au passage \ill{} l'inclusion \ill{}

\ill{} la construction \ill{} $\Sigma$ en \ill{} que \ill{} l'ens. ordonné $\mathcal{M}$ \add{et de la relation $R$ de compatibilité} (sans connaître \add{toutes} les parties $\Phi$ de $\mathcal{M}$, correspondant aux figures). En effet, en termes de $(\mathcal{M}, \leq, R)$, le cosupport d'une partie $\mathcal{M}'$ de $\mathcal{M}$ est défini par
\[
  \operatorname{cosupp}_{\mathcal{M}}(\mathcal{M}') = \lbrace X \in \mathcal{M} \mid (X, X') \in R \ \forall X' \in \mathcal{M}' \rbrace
\]
\note{la longue note oblique de la marge gauche, qui court sur toute la hauteur de la page, n'est lue qu'en fragments (« \ill{} ensembliste \ill{} les exemples \ill{} $\Sigma$ \ill{} », avec les lettres $P$, $Q$, $R$) ; la relation écrite dans la formule est bien $R$, la compatibilité, et non la disjonction}

\emph{Ex} des supports : parties saturées de $\mathcal{M}$ : \ill{} comp. conn. de $\mathcal{M}$ (utiliser C4).
\note{écrit en oblique au bas de la page, séparé du reste par un trait}

\page{69}
J'en viens à la rel. d'ordre $\preccurlyeq$ (subdivision) de nature très différente. Il n'y a pas de plus petit (ni de plus grand) élément, pas de Sup de deux \struck{\ill{}} \add{$F'$, $F''$} \struck{$\preccurlyeq$} toujours (i.e. subdivisions d'une même figure $F$), mais parfois un inf.
\note{le signe de la subdivision, un $\leq$ dont les traits sont incurvés, est rendu $\preccurlyeq$ dans toute la suite}

On dit que deux subdivisions $F'$, $F''$ de $F$ sont \emph{\uncertain{compatibles}} \struck{quand et que seulement} (je ne dirai pas « compatibles », car le terme est déjà pris) si elles admettent une subdivision commune.
\note{le mot souligné est surchargé, et la parenthèse qui suit dit qu'il ne s'agit pas de « compatibles » : la lecture est à reprendre sur la page}
\marginal{On sera amené à \ill{} \ill{} \ill{} les contrées modérées ! \quad $F''' \preccurlyeq F' \preccurlyeq F$, \ $F''' \preccurlyeq F'' \preccurlyeq$}
\note{note oblique de la marge gauche, avec un petit schéma de subdivisions emboîtées dont le second terme de la seconde ligne n'est pas écrit}

On dit que la contrée est \emph{modérée} si deux \add{subdivisions} \struck{\ill{}} de $F$ sont toujours \uncertain{comparables}. \struck{Dans le cas général,}

\struck{\ill{}} Je ne saurais donner \struck{\ill{}} aucune propriété particulière de $\preccurlyeq$ \struck{tout} seule. Les propriétés que j'ai en vue concernent les relations \add{avec} $\leq$.

\emph{C5} \quad Soit $F' \preccurlyeq F$. Alors pour toute figure $G$, $G$ est disjointe de $F$ ssi elle est disjointe de $F'$. [\emph{NB} Il suffit de le poser pour $G = X \in \mathcal{M}$, à cause de C3]
\marginal{\ill{} après C6, C7, \ill{} des disjonctions}
\note{double trait vertical en marge de C5}

\page{70}
\note{son numéro de page, 70, est écrit en surcharge sur un autre chiffre}
Il revient au même de dire que $F$ et $F'$ ont même cosupport, ou encore, même support. Donc \ill{} un support.

Une autre façon de dire (que $F, G$ disj. $\Longrightarrow F', G$ disj.) est ainsi :

\emph{Corollaire} \quad Le cosupport \struck{de} \add{dans $\mathfrak{F}$} tout objet $G$ de $\mathfrak{F}$ est fermé, non seulement pour $\leq$, mais aussi pour $\preccurlyeq$ --- c'est donc \ill{} \ill{} une \emph{sous}-contrée. Donc itou pour toute partie $\Phi \subset \mathfrak{F}$. En d'autres termes, toute partie de $\mathfrak{F}$ qui est un $\mathfrak{F}$-support (ou « $\mathfrak{F}$-saturée », sous-entendu : en termes de $\leq$, cf p. 67) est automatiquement une sous-contrée.

\struck{un ensemble, et $\Phi$ un ens. de parties de $\mathcal{M}$, satisfaisant à la condition précédente (\ill{}). Alors $\Phi$, muni de la relation d'inclusion satisfait C1, et ses objets irréductibles sont les $A_X$ ($X \in \mathcal{M}$), l'application $X \mapsto A_X$ étant une bijection de $\mathcal{M}$ sur l'ens. des \struck{\ill{}} objets irréductibles de $\Phi$. Dém. $A_X$ est irréductible, car si $A_X = \bigcup_{i \in I} A_i$}
\note{la moitié inférieure de la feuille porte, écrit tête-bêche et séparé du haut par un trait, un début de page numéroté « 55 » (chiffre surchargé), barré de longs traits obliques : c'est le texte donné ici entre les marques de biffure}

\page{71}
Par la suite, il faut faire attention si les axiomes \ill{} passent aux sous-contrées. Introduisons quand même la définition en forme [\ill{}]

\emph{Définition 1} \quad On appelle \emph{sous-contrée} de la contrée $\mathcal{C} = (\mathfrak{F}, \leq, \preccurlyeq)$, \struck{tout} sous-ens. qui satisfait les conditions
\begin{enumerate}
  \item[a)] fermé pour $\leq$, et stable par sup quelc. pour $\leq$ [et à ce titre, s'en correspond, \add{par $\mathfrak{F}' \mapsto \mathfrak{F}' \cap \mathcal{M}$,} aux parties fermées \add{$\mathcal{M}'$} de $\mathcal{M}$, par $\mathcal{M}' \mapsto \mathfrak{F}' = \lbrace F \in \mathfrak{F} \mid F \subset \mathcal{M}' \rbrace$].
  \item[b)] fermé pour $\preccurlyeq$ (ce qui impose une restriction aux parties envisagées $\mathcal{M}'$ de $\mathcal{M}$, \ill{} plus bas) \struck{\ill{} si $X \in F$}
\end{enumerate}
\marginal{il faudra \ill{} des \ill{} (\ill{})}
\marginal{(c) Pour $F' \preccurlyeq F$ subd. fixée, l'application $G \mapsto G'$, sous-fig. de $F$ $\to$ sous-fig. de $F'$, \ill{} aux Sup (pour $\leq$) i.e.}
\note{au-dessus de la biffure de b) sont insérés deux mots, « qui ne » et un mot souligné non lu ; le « (c) », entouré, est en marge droite, d'où une flèche descend vers l'énoncé C6}

\emph{C6} \quad Soit $F' \preccurlyeq F$, $G \leq F$ dans $\mathfrak{F}$.

(a) Alors il existe un unique $H \in \mathfrak{F}$ \ill{} :
\[
  F' \preccurlyeq F , \quad G' \preccurlyeq G , \qquad G' \leq F' , \quad G \leq F
\]
(b) et les strates de $G'$ sont les strates $X'$ de $F'$ telles que $X' \ll G$.
\note{les lettres (a) et (b) sont entourées ; sur la page, les relations sont disposées en carré, $F' \preccurlyeq F$ au-dessus de $G' \preccurlyeq G$, avec des signes $\leq$ verticaux ; l'énoncé nomme $H$ la figure que le carré appelle $G'$ ; (b) est bordé d'un double trait, et sous « $X' \ll G$ » une ligne ajoutée n'est lue qu'en partie (« \ill{} complètement \ill{} ») ; un trait vertical borde C6 à gauche}
\marginal{\ill{} symétrie \ill{} $\leq$, $\preccurlyeq$ \ill{} sur la sous-figure $G$ de $F$}

On dit que $G'$ est la \emph{subd.\ de $G$ induite} par la subd. $F'$ de $F$, \ill{} que $G'$ est la \struck{sous-fig.}

\page{72}
sous-figure de \struck{\ill{}} la subdivision $F'$ de $F$, induite par la sous-figure $G$ de $F$.
\marginal{\ill{} \ill{} \ill{}}
\note{note oblique de la marge gauche, en haut de page, non lue}

\emph{C$7_0$} \quad Pour toute figure $F$, sous-figure $G$, et subdivision $G'$ de $G$, il existe une subdivision $F'$ de $F$ qui induise $G'$.

(mais l'énoncé symétrique, où on échangerait les rôles de $\leq$ et $\preccurlyeq$, est faux : il n'est pas vrai, pour $F' \preccurlyeq F$ donnés, que ttes sous-figures $G'$ de $F'$ soient induites par une ss-figure $G$ de $F$).

On \struck{peut} renforcer \uncertain{un peu} \struck{\ill{}} C$7_0$ en C7, en prenant un choix canonique de $F'$.

\struck{\ill{}} Soient $F$, $G$ deux figures. Considérons les conditions \struck{équivalentes} :
\begin{enumerate}
  \item[a)] $G$ est une sous-figure d'une subd. de $F$ : $G \leq F' \preccurlyeq F$
  \item[b)] $G$ est une subdivision d'une sous-figure de $F$ : $G \preccurlyeq G' \leq F$.
\end{enumerate}
On a alors b) $\Rightarrow$ a) (Mais pas nécess. a $\Rightarrow$ b !!!)

\struck{a) $\Rightarrow$ b) résulte de C6}

b) $\Rightarrow$ a) résulte de C$7_0$.
\note{ce dernier alinéa, depuis « Soient $F$, $G$ », est traversé de deux longs traits obliques ; en marge droite un petit schéma : $F$ et $F'$ reliés par un signe de subdivision barré, $G \leq F$ au-dessous, et $G'$ sous $F'$ relié par des pointillés à $G$ ; le mot barré qui précède « Soient » est illisible (on devine « Corollaire »), et le mot souligné ajouté dessous est « considérons »}

\page{73}
\struck{Corollaire 2} \struck{\ill{}}

\emph{Définition} \quad On dit \struck{\ill{}} que $G$ \emph{raffine} $F$ ou est un \emph{raffinement} de $F$ \struck{. On écrit \ill{}} si $\exists\, F'$ avec \struck{$G \ll F$} $G \leq F' \preccurlyeq F$. On écrit $G \ll F$.
\note{au-dessus de « raffine $F$ » une insertion biffée n'est pas lue ; il écrit le signe du raffinement à deux chevrons, parfois à trois (ici même, et page 73 dans « Transitivité ») ; il est rendu partout $\ll$, comme au lot 5}

\emph{Corollaire} \struck{\ill{}} \quad La relation $G \ll F$ est une relation d'ordre (impliquée par $G \leq F$ et par $G \preccurlyeq F$). Utilise C6, ni C5 ni C7.

Transitivité : Soit $F \ll G \ll H$, donc
\[
  \begin{array}{ccc}
    H & \succcurlyeq & H' \\
    & & \geq \\
    G & \succcurlyeq & G' \\
    \geq & & \\
    F & &
  \end{array}
  \qquad \text{on complète par C6} \qquad
  \begin{array}{ccccc}
    H & \succcurlyeq & H' & & \\
    \geq & & \geq & & \\
    \overline{H} & \succcurlyeq & G & \succcurlyeq & G' \\
    \geq & & \geq & & \geq \\
    K & \succcurlyeq & \overline{G} & \succcurlyeq & F
  \end{array}
  \quad (*)
\]
\note{les deux schémas sont des carrés de relations, les signes $\geq$ étant écrits verticalement ; le premier est lu d'après sa disposition (« $H \succcurlyeq H'$ », « $G \succcurlyeq G'$ », $F$ sous $G'$), le second a des points aux sommets complétés ; au-dessus, une première esquisse du même schéma est biffée de longs traits obliques}

et \ill{}, on a $F \preccurlyeq K \leq H$.

Supposons maintenant que $F = H$ i.e. $F \ll G$ et $G \ll F$, prouvons $F = G$. On trouve
\[
  F \preccurlyeq K \leq \underset{\overset{\shortparallel}{H}}{F}
\]
on va en déduire $K = F$ et comme $K \leq \overline{H} \leq F$, il en résultera $\overline{H} = F$ \struck{et donc \ill{} (comme $F \preccurlyeq \overline{G} \leq K$) $\overline{G} = F$} donc $G \preccurlyeq F$. De façon symétrique on aura $F \preccurlyeq G$, d'où $F = G$. Donc il reste à prouver le

\emph{Lemme} \quad Soient $F, K \in \mathfrak{F}$ avec $F \preccurlyeq K \leq F$. Alors $K = F$ (i.e. si $F$ est une subdivision d'une sous-figure $K$, alors $F = K$)

\page{74}
Cela résulte du Corollaire qui suit et de C5

\emph{Corollaire} (Propriété de la relation $\leq$ !) \quad Soient $F$ une figure, $K$ une sous-figure \struck{\ill{}} telle que $K \neq F$. Alors il existe une \struck{\ill{}} figure $G$ (une multistrate, si on veut) qui est \struck{donc} disjointe de $K$, mais non de $F$. \struck{\ill{}} En d'autres termes, $\operatorname{supp} K \neq \operatorname{supp} F$.

Ça n'est pas « formel » en termes \add{des axiomes C1 à C4} sur $(\mathfrak{F}, \leq)$. Supposons p.ex. que $\mathfrak{F}$ admette un plus grand élément, donc il est défini comme l'ens. de \emph{toutes} les parties fermées de l'ens. ordonné $\mathcal{M}$, et « disjoint » signifie ici $A \cap B = \emptyset$ sans plus (la relation de compatibilité $R$ est la relation tjrs vraie). Mais si $F$ est une partie fermée, $K \subset F$ une sous-partie fermée \add{$\neq F$}, il n'est pas clair qu'on puisse trouver une partie fermée $G$ telle que $G \cap K = \emptyset$, $G \cap F \neq \emptyset$. \struck{Par ex.} Si $\mathcal{M}$ a un plus petit élément $0$, \struck{p.ex. $\mathcal{M} = [0, N]$, $R$ avec l'ordre habituel, $F = \mathcal{M}$, $K = [0, 1]$,} Alors pour \emph{toute} partie fermée non vide, les $G$ tels que $F$ et $G$ disjoints sont les $G$ vides, donc ça ne change pas si on remplace $F$ par $K \in F$, \struck{\ill{}} \uncertain{pourvu} que $K$ n'est pas vide

\page{75}
Dans ce cas, l'axiome C6 (pour une relation \add{d'ordre} quelconque $\preccurlyeq$ sur $\mathfrak{F} = \mathfrak{P}_f(\mathcal{M})$) signifie seulement que si $F' \preccurlyeq F$, $F'$ est vide ssi $F$ l'est. Prenons p.ex. \struck{\ill{}} $\preccurlyeq$ \struck{égal à} la relation d'ordre discrète, ça OK. \struck{C$7_0$ est tautologique} C6, C7 tautologiques. On voit donc qu'il manque un axiome, pour assurer la validité du corollaire. On désire que ce soit la \ill{} de Cont 4 (p. 35), \add{\ill{}} Cont 7 : \ill{} pour l'instant \ill{} sais pas utiliser que la relation $\ll$ soit \ill{} \ill{} la propriété suivante, qu'il faudrait \ill{} \ill{} \ill{} :
\marginal{Aussi cela \ill{} résulte \ill{} de C4 \ill{} (\ill{} p. 47)}
\note{« Cont 4 (p. 35) » renvoie au lot précédent (page 35) ; la note de la marge gauche est lue en partie seulement}

\struck{Soient $Y, X \in \mathcal{M}$, si $Y \ll X$, et supposons que $\nexists\, Y' \leq Y$, $Y' \in \partial X$, $Y \cap \partial X =$ \ill{} \ill{} $Y' \leq Y$, $Y' \ll \partial X$ \ill{} les $Y' \leq Y$ ($Y' \in \mathcal{M}$) tels que \ill{} Soient $F$ figure, $F'$ sous-figure, $G \ll F$, $G'$ \add{$\leq G$} le raffinement de $F'$ induit par $G$ (cf plus bas). \struck{supposons,} Alors $G' = \emptyset$ \struck{Alors} ssi $G$ disjoint de $F'$. \struck{$G$ \ill{}}}
\note{tout ce passage est barré de longs traits obliques ; une partie en est encadrée à gauche}

Il faut développer la notion de \emph{raffinement induit}. J'ai envie de \struck{\ill{}} \add{\ill{}} l'énoncé (quitte à rajouter des axiomes) :

\page{76}
\struck{\ill{}} \emph{Prop.} \quad Soit
\[
  \begin{array}{ccc}
    G & \ll & F \\
    & & \geq \\
    & & F'
  \end{array}
\]
alors il existe une plus grande sous-figure $F'$ de $G$, qui raffine $F'$ (et qu'on appelle \emph{raffinement de $F'$ induit} \struck{\ill{}} par le raffinement $G$ de $F$, ou la sous-figure de \struck{$G$ induite} raff. $G$ de $F$, induite par la sous-figure $F'$ de $F$). \struck{\ill{}} Les strates \struck{\ill{} tautologiques} sont les strates de $G$ qui raffinent $F'$. \struck{De plus, $G$ et $F$}
\note{sous $G$, dans le premier schéma, un $\geq$ et un $G'$ sont biffés ; « une plus grande sous-figure $F'$ de $G$ » : lire $G'$, comme le schéma suivant}
\[
  \begin{array}{ccc}
    G & \ll & F \\
    \geq & & \geq \\
    G' & \ll & F'
  \end{array}
\]
De plus, $G'$ est un Inf de $G$ et $F'$ pour la relation $\ll$, en d'autres termes, si \struck{$G$ est} $K$ est une figure telle que $K \ll G$ et $K \ll F'$, alors $K \ll G'$.
\marginal{Dire (on posera) que $G' = \emptyset \Longleftrightarrow G$ disjoint de $F'$ (seul $\Longrightarrow$ non évident…) \struck{$\exists\, G$ telle que}}

\struck{Dém.} Soit $H$ une sous-figure \add{\ill{}} telle que $H \ll F'$, donc pour \uncertain{ttes} strates $X \leq H$, on a $X \ll F'$, donc si on désigne par $G'$ la sous-figure de $G$ engendrée par les strates $Y \leq G$ telles que $Y \ll H$, alors \struck{\ill{}} $H \leq G$ puisque
\[
  H = \operatorname*{Sup}_{\substack{X \in \mathcal{M} \\ X \leq H}} X \leq \operatorname*{Sup}_{\substack{X \in \mathcal{M} \\ X \in G'}} X = G' .
\]
D'autre part, $G' \ll$
\note{la démonstration est barrée de longs traits obliques ; « $Y \ll H$ » et « $H \leq G$ » sont tels sur la page, là où l'on attend $F'$ et $G'$}

\page{77}
On a par définition de $\ll$
\[
  G \leq \overline{F} \preccurlyeq F
\]
\marginal{$G$ sous-figure d'une subd. $\overline{F}$ de $F$}
et on va appliquer C6
\[
  \begin{array}{ccccc}
    G & \leq & \overline{F} & \preccurlyeq & F \\
    \geq & & \geq & & \geq \\
    G' & \leq & \overline{F}' & \preccurlyeq & F' \\
    \shortparallel & & & & \\
    G \cap \overline{F}' & & & &
  \end{array}
\]

\struck{\ill{}} Prouvons que $G'$ est un inf pour $\ll$.

Soit $K$ telle que $K \ll G$, $K \ll F'$. On doit prouver $K \ll G'$. Pour ceci, on doit prouver d'abord $K \ll \overline{F}'$, et on est ramené aux deux cas particuliers où soit $G \preccurlyeq F$, soit $G \leq F$. \ill{}

\emph{Lemme 1} \quad Soit
\[
  \begin{array}{ccc}
    G & \preccurlyeq & F \\
    \geq & & \geq \\
    G' & \preccurlyeq & F'
  \end{array}
  \qquad (\text{cf C6})
\]
et soit $K$ tel que $K \ll G$ et $K \ll F'$, alors $K \ll G'$.

En effet, par \struck{déf.} \add{\ill{}} on a $K \leq \overline{G} \preccurlyeq G$
\[
  \begin{array}{ccccccc}
    K & \leq & \overline{G} & \preccurlyeq & G & \preccurlyeq & F \\
    & & \geq & & \geq & & \geq \\
    & & \overline{G}' & \preccurlyeq & G' & \preccurlyeq & F'
  \end{array}
\]
pour prouver $K \ll G'$, il suffit de prouver $K \leq \overline{G}'$. \struck{\ill{} $X \in \mathcal{M}$, $\forall X \in G$} \ill{} que pour $X \in \mathcal{M}$, $X \leq K \Longrightarrow X \leq \overline{G}'$, or cela résulte de C6, (b) appliqué à :
\[
  \begin{array}{ccc}
    \overline{G} & \preccurlyeq & F \\
    \geq & & \geq \\
    \overline{G}' & \preccurlyeq & F'
  \end{array}
\]
\note{sous « $X \leq K$ » il ajoute : « donc $X \leq \overline{G}$, $X \ll F'$ » ; dans les schémas de cette page les signes $\geq$ sont écrits verticalement entre les lignes}

\page{78}
\emph{Lemme 2} \quad Soit
\[
  \begin{array}{ccc}
    G & \leq & F \\
    \geq & & \geq \\
    G' & \leq & F' \\
    \shortparallel & & \\
    G \cap F' & &
  \end{array}
\]
et soit $K$ tel que $K \ll G$ et $K \ll F'$, alors $K \ll G'$.

On écrit encore $K \leq \overline{G} \preccurlyeq G$, et …

\uncertain{J'abandonne ici} ! Il \ill{} mieux se reporter \ill{} page 51.

\marginal{\emph{NB} \quad $F' \mapsto G'$, sous-fig $F$ $\to$ sous-fig $G$, commute aux Sup$_{\leq}$ quelc., \ill{} et aux Infs$_{\leq}$ quelconques…}
\note{le NB, souligné deux fois, est en marge droite, séparé du texte par un triple trait vertical}

\struck{Soit $F$ une figure}

\struck{\emph{C 8} \quad Soit $F$ une figure, et $X$ une multistrate telle que $X \ll F$ (i.e. $X \leq F' \preccurlyeq F$). Alors $\exists\, Y$ \struck{\ill{}} de $F$, telle que $X \ll Y$.}
\note{C8 est barré de longs traits obliques ; sous « C 8 » un mot biffé n'est pas lu}

\emph{Proposition} \quad Soient $F$ une figure, $X \in \mathcal{M}$, tels que $X \ll F$.
\begin{enumerate}
  \item[a)] \struck{Alors, \ill{}} $\exists\, Y \in \mathcal{M}$, $Y \leq F$ tel que $X \ll Y$
  \item[b)] Parmi les $Y$ satisfaisant cette condition, il y \struck{\ill{}} a un plus petit élément, \struck{\ill{}} \struck{caractérisé par} caractérisé par la propriété que la sous-figure de $X$ induite par $Y$ est $\neq X$.
\end{enumerate}

\emph{Démonstration} \quad $F = \operatorname{Sup}_{Y_i \in \overline{F}} Y_i$, donc \struck{$X$} (par C6 (c)) $X = \operatorname{Sup}_i G_i$ où $G_i$ est la ss-fig. de $X$ induite par les $Y_i$, donc $\exists\, i$ tel que $X \in G_i$, donc $X \ll Y_i$.

\page{79}
Parmi les $Y$ \add{$\in \overline{F}$} tels que $X \ll Y$, choisissons un élément minimal (cela existe, si on admet que la dim. combinatoire des figures irréd. (i.e. multistrates) est finie --- on y reviendra \ill{} \ill{} plus tard encore…) \struck{\ill{}}

Dire que $Y$ est minimal signifie que pour $Y'$ \struck{multistrate} \add{\ill{}} de $Y$, on n'a pas $X \ll Y'$ i.e. $X \neq F_{Y'}$, sous-figure de $X$ induite par $Y'$. Mais cela signifie \ill{} que
\[
  \partial Y = \operatorname*{Sup}_{\substack{Y' \in \mathcal{M} \\ Y' \leq Y}} Y'
\]
est \ill{} par l'image \ill{} de $\partial Y$ différent de $X$ (appliquer C6 (c)).
\note{sous le Sup, une première condition sur $Y'$ est biffée}

\struck{Si \ill{}} Si maintenant $Y' \in \overline{F}$ tel que $X \ll Y'$, prouvons que \struck{$X$} $Y \leq Y'$. \struck{On} En effet l'image inverse $F_{Y'} = X$, \struck{\ill{}} or
\[
  \underset{\overset{\shortparallel}{X}}{F_Y} \cap \underset{\overset{\shortparallel}{X}}{F_{Y'}} = F_{Y \cap Y'}
\]
à cause des propriétés « fonctorielles » des Inf pour $\ll$, donc $F_{Y \cap Y'} = X$, i.e. $X$ se factorise par $Y \cap Y'$ donc par un $Z \leq Y \cap Y' \leq Y$, et donc $Z = Y$ (condition minimale) et par suite $Y = Y'$, \ill{} !

\page{80}
\emph{Corollaire} \quad Soient $F$, $G$ avec $F \ll G$. Pour tout $X \in \overline{F}$, soit $\varphi(X) \in \overline{G}$ la plus petite \struck{multi}strate \add{$Y$} de $G$ telle que \struck{$\varphi(X) \ll$} $X \ll Y$. L'application
\[
  \varphi : \overline{F} \to \overline{G}
\]
ainsi obtenue est \emph{croissante} (trivial).

Jusqu'à maintenant, les axiomes \struck{\ill{}} C5, C6, C7 \ill{} pour $\preccurlyeq$ sont tautologiques \add{dans \ill{}} $\leq$ \add{si on prend pour $\preccurlyeq$ \ill{} de figures}, l'égalité. Il serait temps d'en venir à des propriétés \struck{telles que} plus spécifiques !
\note{les insertions au-dessus de la ligne sont reliées au texte par une longue accolade et ne se lisent qu'en partie}

\subsection{14 juin}
\note{date de sa main, « 14 juin », soulignée, en marge gauche de l'alinéa qui suit}

\struck{Proposition} \emph{Corollaire de prop. p. 78} \quad Je vais poser $\widetilde{F}$ ald $\overline{F}$ pour l'ens. des $X \in \mathcal{M}$ tels que $X \ll F$ (pour éviter confusion avec notion d'adhérence --- Cf l'\uncertain{ancien} C4 \struck{\ill{}} p. 47)
\note{« ald » : « au lieu de » ; le numéro de page du renvoi « p. 78 » est surchargé ; « Cf l'\uncertain{ancien} C4 p. 47) », souligné, est écrit sous l'insertion, au niveau du titre}

Soient $G' \preccurlyeq G \preccurlyeq F$ alors $\widetilde{G}' \cap \widetilde{F} \subset \widetilde{G}$ (donc $\widetilde{G}' \cap \widetilde{F} = \widetilde{G}' \cap \widetilde{G}$).\note{les deux signes de $G' \preccurlyeq G \preccurlyeq F$ sont repassés à l'encre ; lus comme le signe de la subdivision, comme au lot 5, mais le premier pourrait être un $\leq$ transformé} En d'autres termes, toute \add{(multi)} strate de $X$ \struck{\ill{}} qui raffine la sous-figure $G$ de $F$ \add{\ill{}} \add{(multi)} strates de $F$ est strate de $G$.

\struck{Soit $X$ strate de $G'$ qui est strate de $F$}

Dém. \quad $X$ est une strate d'un raffinement $\overline{F}$ de $F$ \add{\ill{}} induit par un raffinement $\overline{G}$ de $G$,
\[
  \begin{array}{ccccc}
    G & \leq & F & \geq & X \\
    \succcurlyeq & & \succcurlyeq & & \\
    \overline{G} & \leq & \overline{F} & & \\
    \geq & & & & \\
    X & & & &
  \end{array}
\]
Supposons \uncertain{$X \leq F$}, \struck{considérons les images inverses} l'image inverse de
\marginal{Si, si \quad C(C\ill{})}
\note{la démonstration est traversée de trois longs traits obliques ; les signes $\succcurlyeq$ du schéma sont écrits verticalement}

\end{document}
